\documentclass[11pt]{article}

% ====================================================
% PACKAGES
% ====================================================
\usepackage[margin=1in]{geometry}
\usepackage[T1]{fontenc}
\usepackage{lmodern}
\usepackage{microtype}
\usepackage{setspace}
\usepackage{parskip}

\usepackage{amsmath, amssymb, mathtools}
\usepackage{siunitx}
\usepackage{graphicx}
\usepackage{float}
\usepackage[justification=centering]{caption}

\usepackage{enumitem}
\usepackage{booktabs}

\usepackage{xcolor}
\usepackage{hyperref}

\usepackage{fancyhdr}
\usepackage{lastpage}
\usepackage{indentfirst}

% ====================================================
% Assignment Info
% ====================================================
\newcommand{\Name}{Arael Anaya}
\newcommand{\Assignment}{Friction Force to Coefficient of Friction}
\newcommand{\course}{Senior Design}

% ====================================================
% MATH SHORTHAND
% ====================================================
\newcommand{\dth}{\dot{\theta}}
\newcommand{\tauf}{\tau_f}
\newcommand{\Nc}{N_c}      % axial clamp load from the nut
\newcommand{\Rp}{R_p}      % radial pin reaction

% ====================================================
% PAGE STYLE
% ====================================================
\onehalfspacing%
\setlength{\parindent}{2em}

\pagestyle{fancy}
\fancyhf{}
\lhead{\Assignment\ |\ \course}
\rhead{\Name\ |\ \today}
\cfoot{Page \thepage\ of~\pageref{LastPage}}

\hypersetup{
    colorlinks,
    linkcolor={red!50!black},
    citecolor={blue!50!black},
    urlcolor={blue!80!black}
}

% ====================================================
% TITLE
% ====================================================
\title{Recovering the Hinge Coefficient of Friction\\
\large from a Measured Total Friction Force}
\author{\Name\\ \course}
\date{\today}

\begin{document}
\maketitle

\noindent
Notation follows the scissor-boom derivation: $\theta$ is the fold angle (the single generalized
coordinate), $N_s$ is the number of hinge-spring joints, and $\tauf$ is the per-joint friction
torque. The equation of motion treats friction only as a lumped constant torque,
$Q_F = N_s\tauf$, so the model never defines a normal force. This note builds the missing link: it
converts a measured total friction force into a per-joint torque and then into a coefficient of
friction $\mu$.

\section{Why the Normal Force Is Not Obvious}

The Lagrangian contains only moments at the joints: the hinge-spring moment $k\theta$, the kink
moment $M_k$, and friction $\tauf$. There is no gravity and no external force on orbit. A pure
moment transmits no force through a pin, so in the idealized model the pin carries essentially no
radial reaction. The normal force that produces friction therefore comes from two sources:
\begin{enumerate}[label=(\roman*)]
    \item \textbf{Fastener clamp} $\Nc$: the preload from the locknut, loaded through the nylon
    washers and the shoulder screw. This is present even with no applied load.
    \item \textbf{Radial pin reaction} $\Rp$: from real loads (weight in the 1-g test rig, tube end
    forces, neighbouring link forces). It is zero in the on-orbit idealization.
\end{enumerate}

\section{From Total Friction to Per-Joint Torque}

\subsection*{Measured quantity}

A force gauge reading $F$ acting perpendicular to the moving panel at lever arm $\ell$ from the
hinge gives the measured total friction generalized torque about $\theta$:
\begin{equation}
Q_{F,\text{meas}} = F\,\ell
\end{equation}

\subsection*{Virtual work over the joints}

Joint $j$ has relative rotation $\varphi_j = c_j\,\theta$. For a scissor, an end pin rotates by
$\theta$ ($c_j = 1$) and a center pin by $2\theta$ ($c_j = 2$). With friction opposing $\dth<0$,
the virtual work of the joint friction torques is
\begin{equation}
\delta W_F = \sum_j \tau_j\, \delta\varphi_j
= \Big(\sum_j c_j\,\tau_j\Big)\delta\theta
\quad\Longrightarrow\quad
Q_F = \sum_j c_j\,\tau_j
\label{eq:QFsum}
\end{equation}
Assuming every joint has the same friction torque $\tauf$:
\begin{equation}
\boxed{\tauf = \frac{Q_{F,\text{meas}}}{\sum_j c_j}}
\label{eq:tauf}
\end{equation}
The existing model sets $c_j=1$ for all $N_s$ joints, so $\sum_j c_j = N_s$ and
$\tauf = Q_F/N_s$. If center pins are counted with $c_j=2$, use the sum in \eqref{eq:tauf}.

\section{Friction Torque Model for One Joint}

\subsection*{Pin (journal) contribution}

The friction torque on a pin of radius $r_p$ carrying radial reaction $\Rp$ is
\begin{equation}
\tau_{\text{pin}} = \mu\,\Rp\,r_p
\label{eq:pin}
\end{equation}

\subsection*{Washer-face contribution}

The nut clamps the nylon washers axially with load $\Nc$. For uniform contact pressure on an
annulus with inner radius $r_i$ and outer radius $r_o$, the mean friction radius is
\begin{equation}
r_m = \frac{2}{3}\,\frac{r_o^{3}-r_i^{3}}{r_o^{2}-r_i^{2}}
\label{eq:rm}
\end{equation}
and the friction torque is
\begin{equation}
\tau_{\text{washer}} = \mu\,\Nc\,r_m
\label{eq:washer}
\end{equation}

\subsection*{Clamp load from nut torque}

The standard torque--preload relation gives
\begin{equation}
\Nc = \frac{T}{K\,d}
\label{eq:clamp}
\end{equation}
with $T$ the applied nut torque, $d$ the screw diameter, and $K\approx0.2$ the nut factor for dry
steel. $K$ varies by roughly $\pm25\%$, so a load washer is preferred if $\Nc$ is needed
accurately.

\subsection*{Total per-joint friction torque}

\begin{equation}
\tauf = \mu\big(\Rp\,r_p + \Nc\,r_m\big)
\label{eq:taufmodel}
\end{equation}

\section{Solving for $\mu$}

Solving \eqref{eq:taufmodel} for the coefficient:
\begin{equation}
\boxed{\mu = \frac{\tauf}{\Rp\,r_p + \Nc\,r_m}}
\label{eq:mu}
\end{equation}

\subsection*{On-orbit idealization ($\Rp\approx0$)}

\begin{equation}
\mu \approx \frac{\tauf}{\Nc\,r_m} = \frac{\tauf\,K\,d}{T\,r_m}
\end{equation}

\subsection*{1-g test rig}

Obtain $\Rp$ from a free-body diagram of the link at fold angle $\theta$:
\begin{equation}
\Rp(\theta) = \sqrt{R_x^{2}(\theta)+R_y^{2}(\theta)}
\end{equation}
where $R_x$, $R_y$ come from force equilibrium including the weight share carried by that pin.
Because $\Rp$ depends on $\theta$, evaluate \eqref{eq:mu} at each angle and report the spread.

\section{Extracting $\mu$ Without Trusting the Odd $N$}

Hand-tightened hardware gives an unrepeatable $\Nc$. Instead, vary the nut torque $T$ with a
torque driver and measure breakaway torque $\tauf$ on a single joint at each setting. Substituting
\eqref{eq:clamp} into \eqref{eq:taufmodel} gives a line in $T$:
\begin{equation}
\tauf = \underbrace{\frac{\mu\,r_m}{K\,d}}_{\text{slope}}\;T
\;+\;\underbrace{\mu\,\Rp\,r_p + \tau_{0}}_{\text{intercept}}
\label{eq:fit}
\end{equation}
where $\tau_0$ collects drag that does not depend on $T$ (bearing, seal, harness). A linear fit
gives
\begin{equation}
\boxed{\mu = \text{slope}\cdot\frac{K\,d}{r_m}}
\end{equation}
The intercept is not used for $\mu$. A large intercept indicates that $T$-independent drag, not
sliding friction at the washers, dominates the joint.

\section{Feeding the Result Back Into the Model}

The design model uses only the lumped $\tauf$ (\texttt{tau\_friction} in \texttt{Params.py}). Two
corrections follow from this derivation:
\begin{enumerate}[label=(\roman*)]
    \item \texttt{tau\_friction = F\_friction * r\_o} uses the link tube radius $r_o=\SI{50.8}{mm}$
    as the lever arm. Friction acts at the pin or washer radius, so the lever arm should be $r_p$
    (or $r_m$), not $r_o$. Prefer entering the measured $\tauf$ from \eqref{eq:tauf} directly.
    \item The spring and kink sizing,
    \begin{equation}
    k = (\text{spring margin})\frac{\tauf}{\theta_1},\qquad
    M_{k0} = (\text{margin})\frac{N_s\big[\tauf - k\theta_{\min}\big]}{2N_k},
    \end{equation}
    then use the measured value in place of the $\SI{5}{N}$ placeholder.
\end{enumerate}

\end{document}
